\documentclass[12pt]{article}

% ----------------------------------------------------------
% PACKAGES
% ----------------------------------------------------------
\usepackage[a4paper, margin=1in]{geometry}
\usepackage{graphicx}
\usepackage{booktabs}
\usepackage{hyperref}
\usepackage{float}
\usepackage{caption}
\usepackage{subcaption}

% ----------------------------------------------------------
% TITLE
% ----------------------------------------------------------
\title{Bridging the Computational Gap: Indoor Visual Localisation on Resource-Constrained Platforms}

\author{
L. Marewangepo$^{a}$ \quad W.J. Smit$^{b}$ \\
\small $^{a}$Department of Mechanical and Mechatronic Engineering, Stellenbosch University, South Africa \\
\small $^{b}$Department of Mechanical and Mechatronic Engineering, Stellenbosch University, South Africa 
}

\date{}

% ----------------------------------------------------------
% DOCUMENT
% ----------------------------------------------------------
\begin{document}

\maketitle

% ----------------------------------------------------------
% ABSTRACT
% ----------------------------------------------------------
\begin{abstract}
Despite falling hardware costs and proven ROI cases, autonomous mobile robot adoption in industrial environments remains low, with 76\% of companies never having deployed such systems. High upfront costs—particularly for compute-intensive perception systems like visual localisation—are cited as the primary barrier. Indoor visual localisation remains computationally demanding when deployed on resource-constrained platforms such as small mobile robots and low-cost embedded devices. Modern localisation pipelines rely on learned feature extractors and large-scale 3D maps, which provide strong accuracy but impose substantial compute, memory, and latency requirements that typically necessitate expensive GPU-equipped platforms.

This work presents a complete visual localisation system designed for deployment on resource-constrained embedded hardware. Our approach combines multiple optimisation strategies: (1) post-training INT8 quantisation of a learned feature extractor (XFeat) for efficient query-time feature extraction, (2) a custom hierarchical vocabulary tree matcher implemented in C++ for rapid descriptor matching, (3) COLMAP-based sparse 3D map representation for memory efficiency, and (4) cross-precision matching enabling INT8 query descriptors to match against existing full-precision (FP32) map databases. This system-level design maintains backward compatibility with standard mapping workflows while achieving real-time performance on embedded platforms.

We evaluate our complete pipeline on a Raspberry Pi 5 across three TUM RGB-D benchmark datasets. The system achieves 24 FPS localisation with mean accuracy of 8.6 cm and memory footprint under 225 MB on an \$80 embedded platform without GPU acceleration. Comparative analysis against classical SIFT-based approaches demonstrates that our learned feature pipeline achieves superior accuracy and speed on identical hardware. These results demonstrate that careful co-design of feature extraction, map representation, and matching strategies can bridge the gap between state-of-the-art visual localisation algorithms and the constraints of low-cost embedded platforms.
\end{abstract}

\textbf{Keywords:} Visual Localization; Autonomous Mobile Robots; Embedded Systems; XFeat; INT8 Quantization; Vocabulary Trees; COLMAP; Resource-Constrained Platforms

% ----------------------------------------------------------
% 1. INTRODUCTION
% ----------------------------------------------------------
\section{Introduction}

The deployment of autonomous mobile robots (AMRs) in industrial environments faces a critical economic bottleneck. While industrial robot prices have fallen by half over the past decade, adoption remains surprisingly low: recent surveys indicate that 76\% of companies have never deployed an AMR, and only 13\% of warehouses are projected to have deployed at least one by 2030~\cite{WarehouseAutomationStatistics2023,therobotreportInteractAnalysisDowngrades2024}. High upfront costs are consistently cited as the primary barrier, with a typical warehouse robot deployment costing upwards of \$1 million~\cite{alexisWarehouseRobotMomentum2024}. Computational requirements for perception and localisation systems represent a significant fraction of this investment, as robust visual localisation typically requires expensive GPU-equipped compute platforms.

Indoor visual localisation, the problem of determining a robot's 6-DOF pose within a pre-built map using camera observations, is a fundamental enabling technology for autonomous navigation in GPS-denied environments such as factories, warehouses, and office buildings. Modern visual localisation methods leverage learned feature extractors based on convolutional neural networks (CNNs), which offer superior robustness and matching performance compared to classical hand-crafted features. However, these learned approaches typically require substantial computational resources for real-time operation, creating a barrier to deployment on resource-constrained mobile robots where cost, size, weight, and power consumption are critical constraints.

Classical hand-crafted features such as SIFT~\cite{loweDistinctiveImageFeatures2004} provide reliable geometric matching but impose prohibitive processing costs when deployed on embedded platforms without GPU acceleration. Recent learned feature extractors such as SuperPoint~\cite{detoneSuperPointSelfSupervisedInterest2018}, D2-Net~\cite{dusmanuD2NetTrainableCNN2019}, and XFeat~\cite{potjeXFeatAcceleratedFeatures2024} demonstrate improved accuracy and robustness, yet their integration into complete localisation pipelines on embedded hardware remains underexplored. 

This work addresses the practical challenge of deploying state-of-the-art visual localisation on low-cost embedded hardware through system-level co-design. We present a complete pipeline that combines: (1)~a learned feature extractor (XFeat) quantised to INT8 for efficient on-device inference, (2)~a custom hierarchical vocabulary tree implemented in C++ for rapid approximate nearest neighbor matching, (3)~sparse COLMAP-based 3D maps built using full-precision features, and (4)~cross-precision descriptor matching that enables quantised queries to localise against existing full-precision map databases. This approach maintains compatibility with standard SLAM/SfM workflows while achieving real-time performance on embedded platforms.

\subsection{Contributions}

This work makes the following contributions:

\begin{itemize}
\item Complete embedded localisation system: We present an end-to-end visual localisation pipeline optimised for resource-constrained platforms. That achieves real-time performance on a Raspberry Pi superior to classical methods while maintaining competitive accuracy.

\item Cross-precision matching analysis: We demonstrate that INT8-quantised query descriptors can effectively match against full-precision (FP32) map descriptors with minimal accuracy degradation, enabling deployment of quantised models without rebuilding existing map databases.


\item Comprehensive embedded evaluation: We provide detailed performance characterisation including accuracy, frame rate, memory footprint, and latency breakdown across three TUM RGB-D benchmark datasets on a Raspberry Pi~5, establishing concrete performance baselines for embedded visual localisation.

\item Learned vs. classical feature comparison: We demonstrate that quantised learned features (XFeat-INT8) substantially outperform classical hand-crafted features (SIFT) on identical embedded hardware interms of speed and match them in accuracy, challenging the conventional assumption that hand-crafted features are more suitable for resource-constrained deployment.
\end{itemize}

By focusing on the interplay between algorithmic design, numerical precision, data structures, and hardware characteristics, this study provides a framework towards deployable, scalable visual localisation solutions that maintain high accuracy without sacrificing computational feasibility. Our results demonstrate that the computational gap separating state-of-the-art localisation algorithms from practical embedded deployment can be bridged through careful system-level optimisation.


\newpage
% 2. RELATED WORK
% ----------------------------------------------------------
\section{Related Work}


\subsection{Indoor Visual Localization Methods}
\subsection{Local Feature Extraction}

\subsubsection{Hand-Crafted Features}

\subsubsection{Learned Features}

\subsection{Model Optimization for Embedded Deployment}

\subsection{Visual Localization on Resource-Constrained Platforms}
\newpage
% ----------------------------------------------------------
% 3. METHODOLOGY
% ----------------------------------------------------------
\section{Methodology}

\subsection{System Overview}

\subsection{Feature-Based Localization Pipeline}

\subsubsection{Feature Extraction}

\subsubsection{Feature Matching}

\subsubsection{Pose Estimation}

\subsection{Map Construction}

\subsection{Model Optimization Strategy}

\subsubsection{ONNX Export and Verification}

\subsubsection{Post-Training Quantization}

\subsection{Cross-Precision Matching}
\newpage
% ----------------------------------------------------------
% 4. EXPERIMENTAL SETUP
% ----------------------------------------------------------
\section{Experimental Setup}

\subsection{Dataset}

\subsection{Hardware Platforms}

\subsection{Evaluation Metrics}

\subsection{Implementation Details}
\newpage
% ----------------------------------------------------------
% 5. RESULTS
% ----------------------------------------------------------
\section{Results}

\subsection{Model Export Verification}

\subsection{Quantization Impact Analysis}

\subsection{Localization Performance}

\subsubsection{Accuracy Evaluation}

\subsubsection{Computational Performance}

\subsection{Cross-Precision Matching Evaluation}

\subsection{Platform Comparison}


\newpage
% ----------------------------------------------------------
% 6. DISCUSSION
% ----------------------------------------------------------
\section{Discussion}

\subsection{Bridging the Computational Gap}

\subsection{Accuracy-Speed Tradeoffs}

\subsection{Cross-Precision Matching Viability}

\subsection{Practical Deployment Considerations}

\subsection{Limitations}
\newpage
% ----------------------------------------------------------
% 7. CONCLUSION
% ----------------------------------------------------------
\section{Conclusion}

\subsection{Summary}

\subsection{Future Work}
\newpage
% ----------------------------------------------------------
% REFERENCES
% ----------------------------------------------------------
\bibliographystyle{ieeetr}
\bibliography{/home/leroy-marewangepo/OneDrive/Masters/Masters_Stuff/Report/template-stb-masters/bib/Masters.bib}

\end{document}
